% GENERATED from content/data/enrichment.json by mxbuild.enrichment; do not edit.
\mxunit{enrichment-waiting-time-solutions}{Attendre le premier succès}
\mxsubtitle{Terminale · probabilités et suites \textperiodcentered{} Solutions détaillées \textperiodcentered{} 08}

\begin{exercise}
Avec $p=\frac{1}{4}$, calculer $P(T=3)$.
\end{exercise}
\begin{solution}
$(\frac{3}{4})^{2}(\frac{1}{4})=\frac{9}{64}$.
\end{solution}

\begin{exercise}
Avec $p=\frac{1}{4}$, calculer $P(T\leq 3)$.
\end{exercise}
\begin{solution}
$1-P(T>3)=1-(\frac{3}{4})^{3}=\frac{37}{64}$.
\end{solution}

\begin{exercise}
Pour $0<p<1$, calculer $P(T>m+n|T>m)$.
\end{exercise}
\begin{solution}
Le quotient des probabilités de queue est $\frac{(1-p)^{m+n}}{1-p}^{m}=(1-p)^{n}$. Après les m échecs, la loi de l’attente restante est inchangée.
\end{solution}

\begin{exercise}
Une réussite a probabilité $\num{0,2}$ à chaque essai indépendant. Quelle est l’attente moyenne ? Le cinquième essai réussit-il nécessairement ?
\end{exercise}
\begin{solution}
$E(T)=5$ essais. Non : chaque essai a toujours probabilité $\num{0,2}$ de réussir, indépendamment des échecs antérieurs.
\end{solution}

\begin{problem}[Chercher, justifier, conclure]
On lance un dé équilibré jusqu’au premier 6, mais on s’arrête dans tous les cas après 10 lancers. Soit N le nombre de lancers réellement effectués. Calculer $P(N=10)$, la probabilité d’avoir obtenu au moins un 6 et E(N). Pour la moyenne, utiliser $E(N)=\sum  P(N>n)$, de $n=0$ à 9.
\end{problem}
\begin{solution}
$N=10$ si les neuf premiers lancers sont des échecs, quel que soit le dixième : $P(N=10)=(\frac{5}{6})^{9}$. La probabilité d’au moins un 6 vaut $1-(\frac{5}{6})^{10}$. Pour $0\leq n\leq 9$, $P(N>n)=(\frac{5}{6})^{n}$, donc $E(N)=6[1-(\frac{5}{6})^{10}]$, soit environ $\num{5,031}$ lancers. La troncature diminue l’espérance par rapport à 6.
\end{solution}

